\section*{Chapter \ref{ch:b2:frontier-orbitals} --- Frontier Orbitals and Reactivity}

\begin{solution}{exo:b2:frontier-orbitals:1}
Ethene: \omterm{def:b2:frontier-orbitals:frontier}{HOMO} $\alpha + \beta$, \omterm{def:b2:frontier-orbitals:frontier}{LUMO} $\alpha - \beta$. Allyl anion (four electrons): \omterm{def:b2:frontier-orbitals:frontier}{HOMO} $\alpha$ (the
nonbonding level), \omterm{def:b2:frontier-orbitals:frontier}{LUMO} $\alpha - 1.414\beta$. Butadiene: \omterm{def:b2:frontier-orbitals:frontier}{HOMO} $\alpha + 0.618\beta$, \omterm{def:b2:frontier-orbitals:frontier}{LUMO} $\alpha -
0.618\beta$.
\end{solution}

\begin{solution}{exo:b2:frontier-orbitals:2}
Hard: \ce{OH-}, \ce{H+}, \ce{F-}. Soft: \ce{I-}, \ce{RS-}, the carbon of iodomethane.
\end{solution}

\begin{solution}{exo:b2:frontier-orbitals:3}
Butadiene (s-cis conformer accessible) and cyclopentadiene (locked s-cis) react;
penta-1,4-diene is not conjugated and does not. The (2\textit{Z},4\textit{Z}) \omterm{def:b2:frontier-orbitals:diels-alder}{diene} can
reach the s-cis form only by pushing its two inner methyl groups into each other: it
reacts very poorly.
\end{solution}

\begin{solution}{exo:b2:frontier-orbitals:4}
Cyclohex-3-ene-1-carbonitrile: a six-membered ring, double bond between the former C2
and C3 of the \omterm{def:b2:frontier-orbitals:diels-alder}{diene}, the \ce{C#N} group on a carbon next to a ring carbon formed from the
\omterm{def:b2:frontier-orbitals:diels-alder}{dienophile}.
\end{solution}

\begin{solution}{exo:b2:frontier-orbitals:5}
$\Delta\varepsilon = \qty{8}{eV}$: perturbative $2\beta^2/\Delta\varepsilon = \qty{0.25}{eV}$. Exact: lower level
$-6 - \sqrt{16 + 1} = \qty{-10.123}{eV}$, gain $2 \times 0.123 = \qty{0.246}{eV}$.
\end{solution}

\begin{solution}{exo:b2:frontier-orbitals:6}
Iodomethane is a \omterm{def:b2:frontier-orbitals:hard-soft}{soft electrophile}: \omterm{def:b2:frontier-orbitals:control}{orbital control}, attack by the atom with the larger
\omterm{def:b2:frontier-orbitals:frontier}{HOMO} coefficient, carbon. A free carbocation is a \omterm{def:b2:frontier-orbitals:hard-soft}{hard electrophile}: \omterm{def:b2:frontier-orbitals:control}{charge control},
attack in part by nitrogen, which carries more negative charge.
\end{solution}

\begin{solution}{exo:b2:frontier-orbitals:7}
C4 of the \omterm{def:b2:frontier-orbitals:diels-alder}{diene} (largest \omterm{def:b2:frontier-orbitals:frontier}{HOMO} coefficient) bonds to the $\beta$ carbon of propenal (largest
\omterm{def:b2:frontier-orbitals:frontier}{LUMO} coefficient): 2-methoxycyclohex-3-ene-1-carbaldehyde, the two substituents on
neighbouring carbons (``ortho'' product).
\end{solution}

\begin{solution}{exo:b2:frontier-orbitals:8}
A bicyclo[2.2.1]heptene (norbornene) skeleton with the ester group on a carbon of the
former \omterm{def:b2:frontier-orbitals:diels-alder}{dienophile}, pointing towards the new \ce{C=C} bond, under the six-membered ring
(endo), on the side opposite the \ce{CH2} bridge.
\end{solution}

\begin{solution}{exo:b2:frontier-orbitals:9}
The two carbonyl groups lower the \omterm{def:b2:frontier-orbitals:frontier}{LUMO} of the \omterm{def:b2:frontier-orbitals:diels-alder}{dienophile}: its gap with the \omterm{def:b2:frontier-orbitals:frontier}{HOMO} of
cyclopentadiene is much smaller than that of ethene, the stabilisation $\beta^2/\Delta\varepsilon$ of
the transition state is larger, and the barrier lower.
\end{solution}

\begin{solution}{exo:b2:frontier-orbitals:10}
The reaction is stereospecific: the maleate gives the adduct with both ester groups cis
(endo,endo as the major product; it is achiral, a meso compound); the fumarate gives the
trans adduct, one ester endo and one exo, formed as a pair of enantiomers (racemic).
\end{solution}

\begin{solution}{exo:b2:frontier-orbitals:11}
Two $\pi$ bonds become two stronger $\sigma$ bonds: $\Delta_r H^\circ < 0$. Two molecules become one:
$\Delta_r S^\circ < 0$. $\Delta_r G^\circ = \Delta_r H^\circ - T\Delta_r S^\circ$ grows with $T$ and becomes
positive above $T_i = \Delta_r H^\circ/\Delta_r S^\circ$: at high temperature the adduct splits back
(retro-Diels--Alder).
\end{solution}

\begin{solution}{exo:b2:frontier-orbitals:12}
Two filled lone-pair orbitals side by side form a four-electron interaction, which
repels. The molecule turns about the \ce{N-N} bond until the lone pairs are roughly
perpendicular (a gauche conformation), where their overlap, and the repulsion, is
small.
\end{solution}

\begin{solution}{pb:b2:frontier-orbitals:1}
\textbf{1.} A five-membered ring with two conjugated double bonds and a \ce{CH2}: the ring
holds the \omterm{def:b2:frontier-orbitals:diels-alder}{diene} unit in the s-cis conformation.
\textbf{2.} The \omterm{def:b2:frontier-orbitals:diels-alder}{dienophile}, through one of its double bonds.
\textbf{3.} A norbornene skeleton fused to a cyclopentene ring; the two new $\sigma$ bonds join
the ends of the \omterm{def:b2:frontier-orbitals:diels-alder}{diene} to the two carbons of the \omterm{def:b2:frontier-orbitals:diels-alder}{dienophile}'s double bond.
\textbf{4.} The \omterm{def:b2:frontier-orbitals:endo}{endo adduct} (\omterm{def:b2:frontier-orbitals:endo}{endo rule}): the remaining ring of the \omterm{def:b2:frontier-orbitals:diels-alder}{dienophile} lies under
the \omterm{def:b2:frontier-orbitals:diels-alder}{diene} in the transition state.
\textbf{5.} Yes: it is concerted, both bonds form on the same face of each partner, and the
geometry of the \omterm{def:b2:frontier-orbitals:diels-alder}{dienophile} is kept.
\textbf{6.} Cyclopentadiene is both a good \omterm{def:b2:frontier-orbitals:diels-alder}{diene} (locked s-cis) and a \omterm{def:b2:frontier-orbitals:diels-alder}{dienophile}, and the
reaction is concerted, with no charged intermediate to stabilise.
\textbf{7.} Negative: two $\pi$ bonds are replaced by two $\sigma$ bonds.
\textbf{8.} Negative: two molecules give one.
\textbf{9.} $\Delta_r G^\circ = \Delta_r H^\circ - T\Delta_r S^\circ$ with both terms negative: negative at low
$T$, zero at $T_i = \Delta_r H^\circ/\Delta_r S^\circ$, positive above.
\textbf{10.} At room temperature $T < T_i$: dimerisation is favoured, though slow. Near
the boiling point of the dimer the equilibrium favours the monomer, and it is reached
quickly.
\textbf{11.} The monomer leaves the hot mixture as vapour: removing a product keeps the
equilibrium shifting towards the monomer (Le Chatelier).
\textbf{12.} At room temperature it dimerises again; cold slows that reaction.
\textbf{13.} $0.62 - (-0.62) = 1.24|\beta|$.
\textbf{14.} $0.62 - (-0.35) = 0.97|\beta|$.
\textbf{15.} Cyclopentadiene with maleic anhydride: the smaller gap gives a larger
stabilisation of the transition state, a faster reaction.
\textbf{16.} They are conjugated with the \ce{C=C} bond and electron-withdrawing: as in
propenal, they draw the $\pi^*$ level down.
\textbf{17.} Its \omterm{def:b2:frontier-orbitals:frontier}{LUMO}, which has lobes on the carbonyl carbons facing the central carbons of
the \omterm{def:b2:frontier-orbitals:diels-alder}{diene}'s \omterm{def:b2:frontier-orbitals:frontier}{HOMO}.
\textbf{18.} Norbornene skeleton; the anhydride ring under the new double bond, opposite the
\ce{CH2} bridge.
\textbf{19.} The two bridgehead carbons and the two carbons carrying the anhydride: four
stereocentres.
\textbf{20.} No: a mirror plane passes through the \ce{CH2} bridge and between the two halves
of the molecule; it is a meso compound.
\textbf{21.} They were cis on the \omterm{def:b2:frontier-orbitals:diels-alder}{dienophile}, and the concerted addition keeps them cis.
\textbf{22.} The same skeleton with the anhydride ring on the side of the \ce{CH2} bridge.
\textbf{23.} No: changing endo to exo would mean breaking and remaking \ce{C-C} bonds, through
the retro-reaction, which needs heating.
\textbf{24.} Water opens the anhydride: the cis dicarboxylic acid, both \ce{COOH} groups endo.
\textbf{25.} The \omterm{def:b2:frontier-orbitals:endo}{endo adduct} has $\boldsymbol{4}$ stereocentres.
\end{solution}
